\documentclass[12pt]{article}
\usepackage[margin=1in]{geometry}
\usepackage{setspace}
\usepackage{hyperref}
\usepackage{enumitem}
\usepackage{amsmath}
\usepackage{booktabs}
\usepackage{multicol}

\setstretch{1.05}

\title{\textbf{A Machine Learning Approach to Modeling Market Value}}
\author{
\textbf{Name:} Ryan Barton \quad \textbf{NetID:} \texttt{rmb298} \\
\textbf{Team:} Solo
}
\date{}

\begin{document}
\maketitle

%=========================================================
\section*{Research Question and Motivation}
\begin{description}[leftmargin=2.25em]
\item[\textbf{Research Question:}]
Can a ML model predict the fair market valuation of a company from their past performance, and macroeconomic factors?
\item[\textbf{Motivation:}]
Market capitalization is an indicator of a company's financial well‑being.  
Being able to predict this number will give insight into the factors that drive it.
Utilizing open‑source APIs to gather macroeconomic and foundational conditions, we will be able to evaluate how closely these features correlate with market capitalization. 
By using data‑science and machine learning techniques paired with a relational database, we will be able to effectively create a pipeline to compare predicted market capitalization to actual market capitalization.
\end{description}
%=========================================================
\section*{Data Sources}

\noindent\textbf{Planned Data:}
\begin{itemize}[leftmargin=2.25em]
  \item \textbf{Scope:} Current 100 U.S. public companies by market capitalization.
  \item \textbf{Daily Market Data:} Price and volume (OHLCV) for each company. This will be used to calculate momentum/volatility.
  \item \textbf{Quarterly Market Data:} Revenue, net income, assets, liabilities, shares outstanding. Used to create quarterly snapshots
  \item \textbf{Macroeconomic Data:} Unemployement rate, Federal Funds Rate, CPI inflation (YoY). Used for correlation of monthly snapshots
\end{itemize}

\noindent\textbf{How the data will be obtained:}\\
This Data will be obtained through multiple API's. This list will be subject to change if need be for the requirements of this project.
\begin{itemize}[leftmargin=2.25em]
  \item \textbf{Alpha Vantage(API) OR Yahoo Finance(API):} Daily Market Data, Quarterly Market Data.
  \item \textbf{FRED(API):} Macroeconomic Data.
  \item \textbf{100 Companies(Website):} https://companiesmarketcap.com/usa/largest-companies-in-the-usa-by-market-cap/ will be website used to scrape 100 companies for analysis.
\end{itemize}


\noindent\textbf{Cleaning and preprocessing plan:}
\begin{itemize}[leftmargin=2.25em]
  \item \textbf{Normalizing Identifiers:} Ticker symbols are consistent and only map to one company.
  \item \textbf{Fix Types and Dates:} Consistent formats for dates, and cast numeric data to numeric types to save space.
  \item \textbf{Handle insufficent Data:} Handle NaN type, drop columns or rows if necessary. Understand the missing data and how it may skew results.
  \item \textbf{Build Snapshots of Data:} Use as-of methodology to only use data up to quarterly snapshot to avoid dishonest predictions.
\end{itemize}


\noindent\textbf{Potential challenges:}
\begin{itemize}[leftmargin=2.25em]
   \item \textbf{API limits:} Rate limits, and possible missing data will hinder process. May have to narrow scope if need be.
   \item \textbf{Building as-of Snapshots:} This is a complex step that requires understanding the data inorder to prevent ahead of time bias.
\end{itemize}

%=========================================================
\section*{Methodology and Analysis Plan}

\noindent\textbf{Database design (SQL / data management):}\\
\begin{table}[h]
\centering
\caption{SQL Database Schema for Market Valuation Model}
\label{tab:database_schema}
\begin{tabular}{@{}llll@{}}
\toprule
\textbf{Table} & \textbf{Primary Key} & \textbf{Foreign Key} & \textbf{Key Columns} \\ \midrule
\textbf{Companies} & \texttt{ticker} & -- & \texttt{sector}, \texttt{industry\_group} \\
\textbf{Daily\_Prices} & \texttt{(ticker, date)} & \texttt{ticker} & \texttt{close}, \texttt{volume}, \texttt{hl\_pct\_change} \\
\textbf{Fundamentals} & \texttt{(ticker, report\_date)} & \texttt{ticker} & \texttt{net\_income}, \texttt{total\_assets}, \texttt{shares\_out} \\
\textbf{Macro} & \texttt{date} & -- & \texttt{unemployment\_rate}, \texttt{cpi\_yoy}, \texttt{fed\_rate} \\ \bottomrule
\end{tabular}
\end{table}


\noindent\textbf{Feature engineering:}\\
\begin{itemize}
  \item \textbf{P/E ratio:} Price to Earnings ratio computed with daily market values
  \item \textbf{Debt-to-equity:} Can be calculated by quarterly foundamental data.
  \item \textbf{Momentum and Volatility Features:}  utilizing the daily market values we can calculate the standard deviation and the volatility of a stock.
  \item \textbf{Macro Economic Features:}  Determine how the enviroment, i.e. inflation correlates to the market capitalization of a company.
\end{itemize}

\noindent\textbf{Models:}\\
\begin{itemize}
  \item \textbf{Gradient Boosting Machine:} This is the most practical, as it does not require signficant compute power, so we will start with this type of model.
  \item \textbf{Long Short-Term Memory:} This is the less practical option but does offer great insight to some factors such as identifying historical patterns that will affect a company negatively.
  \item \textbf{Plan:} Start with GBM which requires detailed feature engineering, but can handle noiser data. If time permits move into LSTM for market volatility and trends. (7)
\end{itemize}

\noindent\textbf{Tools/libraries:}\\
\begin{itemize}
\begin{multicols}{4}
  \item pandas
  \item numpy
  \item sqlite3
  \item xgboost
  \item yfinance
  \item fredapi
  \item alpha\_vantage
  \item pytorch
\end{multicols}
\end{itemize}

%=========================================================
\section*{Expected:}
The market is notoriously hard to predict. This means that for this project, we will in no way try to achieve 100\% accuracy. 
Instead we will try to strive for a statistically signficant accuracy of 55\%. This shouldn't be unobtainable because the top 100 companies 
are more stable. With these predictions we will be able to see how macroeconmics affect companies values and how hard it is to predict the enviroment. 
For instance, our model will most likely not be able to predict the market crash from COVID-19. However, we do believe that combining macroeconomic and Fundamentals
will increase accuracy of predictions and is a better approach than observing just price history alone. (7)
\section*{Resources:}
\begin{enumerate}
\begin{multicols}{2}
  \item https://fred.stlouisfed.org/
  \item https://www.alphavantage.co
  \item https://pypi.org/project/yfinance/
  \item https://finance.yahoo.com/
  \item https://pandas.pydata.org/
  \item https://pytorch.org/
  \item https://arxiv.org/html/2505.23084v1
  
\end{multicols}
\end{enumerate}


\end{document}